% GENERATED FILE. Edit canonical lessons, then run npm run generate:books.
% canonical-source-hash: fnv1a64:e2dbcfabc3a1d00d
% canonical-lessons: MR-C524-kacra, MR-C524-kersuni, MR-C524-potere, MR-C524-jhatakne, MR-C524-ghasne

\chapter{Cleaning Day}
\label{ch:mr-cleaning-day}
\hlchaptermodality{524}

\begin{chapteropening}
\textbf{By the end of this chapter:} \emph{I can talk about rubbish, sweeping and scrubbing.}

Complete the last lesson of chapter 524: I can talk about rubbish, sweeping and scrubbing.
\end{chapteropening}

\section[\texorpdfstring{\mr{कचरा} (kacrā)}{कचरा (kacrā)}]{\mr{कचरा} (kacrā) — rubbish; litter}

\label{lesson:MR-C524-kacra}



\begin{quote}
Before the new one: say the Marathi for a mat, then the Marathi for to roll up.
\end{quote}

\subsection*{You'll want to know: \mr{कचरा}}
\textbf{\mr{कचरा}} — \emph{kacrā} — \textquotedblleft{}rubbish; litter\textquotedblright{}.

\begin{quote}
\mr{उद्यानात} \mr{खूप} \mr{कचरा} \mr{पडला} \mr{होता}.
\end{quote}

\begin{quote}
\emph{udyānāt khūp kacrā paḍlā hotā.}
\end{quote}

\begin{quote}
There was a lot of litter lying in the park.
\end{quote}

Masculine.

\begin{cousinweb}
\textbf{\mr{कचरा}} (\emph{kacrā}) is a tadbhava: CDIAL lists Marathi \emph{kacrā}, \textquotedblleft{}rubbish\textquotedblright{}, under Sanskrit \emph{kaccara}, \textquotedblleft{}dirty, vile\textquotedblright{} (CDIAL 2615).
\end{cousinweb}

\begin{grammarlens}[title={What you've built}]
One more word for describing what you lived and felt.
\end{grammarlens}

\subsection*{Your turn}
\begin{itemize}
\raggedright
  \item \emph{Say it:} \emph{kacrā}
  \item \emph{Say it:} \emph{kacrā}, once more
  \item \emph{Say it:} \emph{udyānāt khūp kacrā paḍlā hotā.} — There was a lot of litter lying in the park.
  \item \emph{Recall:} say the Marathi for a mat, then the Marathi for to roll up, then say \emph{kacrā} again
\end{itemize}

\begin{tcolorbox}[breakable,colback=teal!4,colframe=teal!35!black,title={Before you move on}]
What does \mr{कचरा} mean? (\textquotedblleft{}Rubbish; litter\textquotedblright{}.) Say it once more.
\end{tcolorbox}
\section[\texorpdfstring{\mr{केरसुणी} (kersuṇī)}{केरसुणी (kersuṇī)}]{\mr{केरसुणी} (kersuṇī) — a grass broom}

\label{lesson:MR-C524-kersuni}



\begin{quote}
Before the new one: say the Marathi for to roll up, then the Marathi for rubbish.
\end{quote}

\subsection*{You'll want to know: \mr{केरसुणी}}
\textbf{\mr{केरसुणी}} — \emph{kersuṇī} — \textquotedblleft{}a grass broom\textquotedblright{}.

\begin{quote}
\mr{आईने} \mr{केरसुणीने} \mr{खोली} \mr{झाडली}.
\end{quote}

\begin{quote}
\emph{āīne kersuṇīne kholī jhāḍlī.}
\end{quote}

\begin{quote}
Mother swept the room with a grass broom.
\end{quote}

Feminine.

Softer than a \mr{झाडू}, for sweeping indoors.

\begin{cousinweb}
\textbf{\mr{केरसुणी}} (\emph{kersuṇī}) is a compound of \mr{केर}, \textquotedblleft{}sweepings\textquotedblright{}, and \mr{सुणी}; the origin of both parts is not settled.
\end{cousinweb}

\begin{grammarlens}[title={What you've built}]
One more word for describing what you lived and felt.
\end{grammarlens}

\subsection*{Your turn}
\begin{itemize}
\raggedright
  \item \emph{Say it:} \emph{kersuṇī}
  \item \emph{Say it:} \emph{kersuṇī}, once more
  \item \emph{Say it:} \emph{āīne kersuṇīne kholī jhāḍlī.} — Mother swept the room with a grass broom.
  \item \emph{Recall:} say the Marathi for to roll up, then the Marathi for rubbish, then say \emph{kersuṇī} again
\end{itemize}

\begin{tcolorbox}[breakable,colback=teal!4,colframe=teal!35!black,title={Before you move on}]
What does \mr{केरसुणी} mean? (\textquotedblleft{}A grass broom\textquotedblright{}.) Say it once more.
\end{tcolorbox}
\section[\texorpdfstring{\mr{पोतेरे} (potere)}{पोतेरे (potere)}]{\mr{पोतेरे} (potere) — a floor cloth}

\label{lesson:MR-C524-potere}



\begin{quote}
Before the new one: say the Marathi for rubbish, then the Marathi for a grass broom.
\end{quote}

\subsection*{You'll want to know: \mr{पोतेरे}}
\textbf{\mr{पोतेरे}} — \emph{potere} — \textquotedblleft{}a floor cloth\textquotedblright{}.

\begin{quote}
\mr{न्हाणीघराबाहेर} \mr{एक} \mr{जुनं} \mr{पोतेरं} \mr{पडलं} \mr{होतं}.
\end{quote}

\begin{quote}
\emph{nhāṇīgharābāher ek junaṁ poteraṁ paḍlaṁ hotaṁ.}
\end{quote}

\begin{quote}
An old floor cloth lay outside the bathroom.
\end{quote}

Neuter.

The dictionary writes \mr{पोतेरे}; a sentence says \mr{पोतेरं}, like \mr{जुनं}.

\begin{cousinweb}
\textbf{\mr{पोतेरे}} (\emph{potere}) is inherited: CDIAL lists Marathi \emph{potrẽ}, \textquotedblleft{}a rag for smearing\textquotedblright{}, under Sanskrit \emph{pōta}, \textquotedblleft{}cloth\textquotedblright{} (CDIAL 8400).
\end{cousinweb}

\begin{grammarlens}[title={What you've built}]
One more word for describing what you lived and felt.
\end{grammarlens}

\subsection*{Your turn}
\begin{itemize}
\raggedright
  \item \emph{Say it:} \emph{potere}
  \item \emph{Say it:} \emph{potere}, once more
  \item \emph{Say it:} \emph{nhāṇīgharābāher ek junaṁ poteraṁ paḍlaṁ hotaṁ.} — An old floor cloth lay outside the bathroom.
  \item \emph{Recall:} say the Marathi for rubbish, then the Marathi for a grass broom, then say \emph{potere} again
\end{itemize}

\begin{tcolorbox}[breakable,colback=teal!4,colframe=teal!35!black,title={Before you move on}]
What does \mr{पोतेरे} mean? (\textquotedblleft{}A floor cloth\textquotedblright{}.) Say it once more.
\end{tcolorbox}
\section[\texorpdfstring{\mr{झटकणे} (jhaṭakṇe)}{झटकणे (jhaṭakṇe)}]{\mr{झटकणे} (jhaṭakṇe) — to shake off; to dust}

\label{lesson:MR-C524-jhatakne}



\begin{quote}
Before the new one: say the Marathi for a grass broom, then the Marathi for a floor cloth.
\end{quote}

\subsection*{You'll want to know: \mr{झटकणे}}
\textbf{\mr{झटकणे}} — \emph{jhaṭakṇe} — \textquotedblleft{}to shake off; to dust\textquotedblright{}.

\begin{quote}
\mr{आजीने} \mr{गच्चीवर} \mr{चादर} \mr{झटकली}.
\end{quote}

\begin{quote}
\emph{ājīne gaccīvar cādar jhaṭaklī.}
\end{quote}

\begin{quote}
Grandmother shook out the sheet on the terrace.
\end{quote}

Past, with -\mr{ने} (he …): \mr{त्याने} \mr{चादर} \mr{झटकली}, \emph{tyāne cādar jhaṭaklī}.

\begin{cousinweb}
\textbf{\mr{झटकणे}} (\emph{jhaṭakṇe}) is inherited: CDIAL sets Marathi \emph{jhaṭakṇẽ} among the words of sudden movement under \emph{*jhaṭ} (CDIAL 5327).
\end{cousinweb}

\begin{grammarlens}[title={What you've built}]
One more word for describing what you lived and felt.
\end{grammarlens}

\subsection*{Your turn}
\begin{itemize}
\raggedright
  \item \emph{Say it:} \emph{jhaṭakṇe}
  \item \emph{Say it:} \emph{jhaṭakṇe}, once more
  \item \emph{Say it:} \emph{ājīne gaccīvar cādar jhaṭaklī.} — Grandmother shook out the sheet on the terrace.
  \item \emph{Recall:} say the Marathi for a grass broom, then the Marathi for a floor cloth, then say \emph{jhaṭakṇe} again
\end{itemize}

\begin{tcolorbox}[breakable,colback=teal!4,colframe=teal!35!black,title={Before you move on}]
What does \mr{झटकणे} mean? (\textquotedblleft{}To shake off; to dust\textquotedblright{}.) Say it once more.
\end{tcolorbox}
\section[\texorpdfstring{\mr{घासणे} (ghāsṇe)}{घासणे (ghāsṇe)}]{\mr{घासणे} (ghāsṇe) — to scrub; to rub}

\label{lesson:MR-C524-ghasne}



\begin{quote}
Before the new one: say the Marathi for a floor cloth, then the Marathi for to shake off.
\end{quote}

\subsection*{You'll want to know: \mr{घासणे}}
\textbf{\mr{घासणे}} — \emph{ghāsṇe} — \textquotedblleft{}to scrub; to rub\textquotedblright{}.

\begin{quote}
\mr{जेवणानंतर} \mr{मी} \mr{सगळी} \mr{भांडी} \mr{घासली}.
\end{quote}

\begin{quote}
\emph{jevṇānantar mī sagḷī bhāṇḍī ghāslī.}
\end{quote}

\begin{quote}
After the meal I scrubbed all the dishes.
\end{quote}

Past, with -\mr{ने} (he …): \mr{त्याने} \mr{भांडी} \mr{घासली}, \emph{tyāne bhāṇḍī ghāslī}.

\begin{cousinweb}
\textbf{\mr{घासणे}} (\emph{ghāsṇe}) is a tadbhava, inherited from Sanskrit \emph{gharṣati}, \textquotedblleft{}rubs\textquotedblright{} (CDIAL 4450).
\end{cousinweb}

\begin{grammarlens}[title={What you've built}]
One more word for describing what you lived and felt.
\end{grammarlens}

\subsection*{Your turn}
\begin{itemize}
\raggedright
  \item \emph{Say it:} \emph{ghāsṇe}
  \item \emph{Say it:} \emph{ghāsṇe}, once more
  \item \emph{Say it:} \emph{jevṇānantar mī sagḷī bhāṇḍī ghāslī.} — After the meal I scrubbed all the dishes.
  \item \emph{Recall:} say the Marathi for a floor cloth, then the Marathi for to shake off, then say \emph{ghāsṇe} again
\end{itemize}

\begin{tcolorbox}[breakable,colback=teal!4,colframe=teal!35!black,title={Before you move on}]
What does \mr{घासणे} mean? (\textquotedblleft{}To scrub; to rub\textquotedblright{}.) Say it once more.
\end{tcolorbox}
